\documentclass{article}
\usepackage[utf8]{inputenc}
\usepackage{amsmath}
\usepackage{graphicx}

\title{示例文档}
\author{测试作者}
\date{\today}

\begin{document}

\maketitle

\section{第一章：介绍}

这是第一章的内容。它包含了一些基本的文本，用于测试各种工具的功能。

\subsection{子章节}

这是子章节的内容。它展示了文档的层级结构。

\section{第二章：主要内容}

这是第二章的内容。它包含了更多的文本和结构。

\subsection{第一个子章节}

这是第一个子章节的内容。

\subsection{第二个子章节}

这是第二个子章节的内容。

\section{第三章：总结}

这是第三章的内容，用于总结文档的主要内容。

\subsection{结论}

这是结论部分的内容。

\section{数学公式示例}

这是一个行内公式：$E = mc^2$

这是一个块级公式：
\begin{equation}
\int_{-\infty}^{\infty} e^{-x^2} dx = \sqrt{\pi}
\end{equation}

\section{列表示例}

\begin{itemize}
  \item 第一项
  \item 第二项
  \item 第三项
\end{itemize}

\begin{enumerate}
  \item 编号项1
  \item 编号项2
  \item 编号项3
\end{enumerate}

\section{表格示例}

\begin{tabular}{|c|c|c|}
\hline
列1 & 列2 & 列3 \\
\hline
数据1 & 数据2 & 数据3 \\
数据4 & 数据5 & 数据6 \\
\hline
\end{tabular}

\section{引用示例}

\begin{quote}
这是一个引用块的示例。
\end{quote}

\section{强调示例}

这是\textbf{粗体文本}，这是\textit{斜体文本}。

\section{链接示例}

这是一个链接示例：\url{https://example.com}

\section{图片示例}

\begin{figure}[h]
\centering
\includegraphics[width=0.5\textwidth]{example.png}
\caption{示例图片}
\end{figure}

\end{document}

